% ==============================================================================
% Section 9: Conclusion
% ==============================================================================

\section{Conclusion}
\label{sec:conclusion}

This paper has presented AR-DeC, a novel mobile application that integrates augmented reality, on-device optical character recognition, adaptive intelligent tutoring, and theoretically-grounded gamification to support vocabulary acquisition during natural reading activities. The system addresses a significant gap in the educational technology literature, where AR, ITS, and gamification have been extensively studied in isolation but rarely combined in a single, mobile-native platform.

The principal contributions of this work are as follows. First, we have proposed and detailed a system architecture that demonstrates the technical feasibility of integrating on-device OCR (Google ML Kit), adaptive tutoring with Bayesian Knowledge Tracing, AR rendering (ViroReact), and gamification within a cross-platform React Native application. Second, we have grounded the system design in established theoretical frameworks---Keller's ARCS motivational model and Bloom's revised taxonomy---providing explicit mappings between theory and system features that can guide future research and development in this area. Third, we have specified a rigorous Design-Based Research evaluation methodology with validated instruments (VKS, SUS, ARCS questionnaire) and a mixed-methods analysis plan capable of addressing questions about vocabulary acquisition, engagement, and adaptive learning effectiveness.

The implications of this work extend beyond the specific AR-DeC implementation. The architectural patterns and design principles described herein---particularly the integration of BKT-driven content adaptation with camera-mediated AR interaction and the application of the ARCS model to gamification design---are transferable to a broad range of mobile educational applications. The on-device processing approach ensures that AR-DeC can function effectively in contexts with limited or no internet connectivity, an important consideration for educational equity.

\subsection{Future Work}
\label{subsec:future_work}

Several directions for future development and research are planned:

\begin{enumerate}
    \item \textbf{Multi-language Support}: Extending AR-DeC to support multiple source and target languages, enabling use by language learners and multilingual students. ML Kit's support for over 100 languages provides a foundation for this extension.

    \item \textbf{Collaborative Features}: Introducing shared vocabulary lists, peer-to-peer challenges, and collaborative quizzes to leverage the social dimension of learning and address the relatedness component of Self-Determination Theory \cite{ryan2000}.

    \item \textbf{AI Model Fine-tuning}: Training specialised language models on domain-specific corpora to improve the quality and relevance of generated explanations, reducing dependence on external API calls and enabling fully on-device content generation.

    \item \textbf{Longitudinal Evaluation}: Conducting extended studies over a full academic semester to assess long-term retention effects, sustained engagement beyond the novelty period, and the impact of AR-DeC on broader academic performance.

    \item \textbf{Accessibility Enhancements}: Incorporating text-to-speech functionality, adjustable font sizes and contrast ratios, and haptic feedback to ensure the application is usable by learners with visual or motor impairments.

    \item \textbf{Advanced Student Modelling}: Exploring deep knowledge tracing models and multi-dimensional student representations that capture not only vocabulary mastery but also reading fluency, topic interest, and optimal learning schedules.
\end{enumerate}

\subsection*{Acknowledgments}

The authors gratefully acknowledge the support of [funding body/institution]. We thank [names] for their contributions to the pilot testing of the application.

\subsection*{Data Availability Statement}

The source code for AR-DeC is available at [repository URL]. Anonymised evaluation data will be made available upon reasonable request following the completion of the study.

\subsection*{Conflict of Interest}

The authors declare no conflict of interest.
